\chapter[Resultado e Discussão]{Resultado e Discussão}

Este capítulo apresenta os resultados obtidos a partir do treinamento e da avaliação dos modelos de redes neurais convolucionais propostos 
neste trabalho. São analisados o desempenho da CNN1D e da CNN2D, considerando as métricas de acurácia, sensibilidade, especificidade, precisão 
e F1-score. Os resultados são apresentados individualmente para cada modelo e, posteriormente, comparados.

\section{Resultados da CNN1D}

Nesta seção são apresentados os resultados obtidos pela CNN1D, aplicada diretamente aos sinais temporais de ECG.

A matriz de confusão obtida no conjunto de teste pode ser visualizada na Figura \ref{fig:matconf1d}.

\begin{figure}[ht]
  \caption{Matriz de confusão da CNN1D}
  \label{fig:matconf1d}
  \centering
  \includegraphics[width=0.7\textwidth]{MatConf1D.jpeg}
  \legend{Elaborado pelo autor (2026)}
\end{figure}

Observa-se que o modelo apresenta bom desempenho geral, com alta concentração de valores na diagonal principal da matriz, indicando boa taxa de acertos 
entre as classes. No entanto, a análise detalhada dos erros revela padrões relevantes de confusão entre classes específicas.

A classe N (batimentos normais) apresentou alguns casos de confusão com as classes S, V e F, o que pode ser explicado pela similaridade morfológica 
entre determinados batimentos, especialmente em situações onde há variações sutis no sinal de ECG. De forma semelhante, a classe V apresentou pequenas 
confusões com a classe F, indicando que o modelo pode ter dificuldade em distinguir padrões mais complexos ou intermediários entre essas categorias.

Um ponto particularmente relevante é a ocorrência de batimentos da classe F (batimentos de fusão) classificados como pertencentes à classe N. Esse tipo 
de erro é criticamente importante do ponto de vista clínico, uma vez que um batimento anômalo sendo interpretado como normal pode levar à subestimação 
de condições cardíacas potencialmente relevantes. Esse resultado sugere que, embora o modelo apresente bom desempenho global, ainda há limitações na 
capacidade de identificar corretamente padrões menos frequentes ou mais complexos.

Por outro lado, a classe Q apresentou desempenho praticamente perfeito, com altíssima sensibilidade e precisão, indicando que seus padrões são mais 
facilmente distinguíveis pelo modelo.

Os valores calculados para cada métrica são apresentados na Tabela \ref{tab:table5}.

\begin{table}[ht]
  \caption{Métricas de desempenho da CNN1D}
  \label{tab:table5}
  \centering
  \resizebox{\textwidth}{!}{%
  \begin{tabular}{cccccc}
    \hline
    \bfseries Classe & \bfseries Acurácia & \bfseries Sensibilidade & \bfseries Especificidade & \bfseries Precisão & \bfseries F1-Score\\
    \hline
    N & 0.9590 & 0.9250 & 0.9675 & 0.8768 & 0.9002\\
    S & 0.9750 & 0.9300 & 0.9862 & 0.9442 & 0.9370\\
    V & 0.9760 & 0.9300 & 0.9875 & 0.9490 & 0.9394\\
    F & 0.9650 & 0.9000 & 0.9812 & 0.9231 & 0.9114\\
    Q & 0.9970 & 0.9950 & 0.9975 & 0.9900 & 0.9925\\
    \hline
    \textbf{MÉDIA} & \textbf{0.9744} & \textbf{0.9360} & \textbf{0.9840} & \textbf{0.9366} & \textbf{0.9361}\\
    \hline
  \end{tabular}%
  }
  \source
\end{table}

De forma geral, o modelo apresentou bom desempenho, com acurácia média de 97,44\% e F1-score médio de 93,61\%. No entanto, a sensibilidade da classe 
F (0,90) é inferior às demais, reforçando a dificuldade do modelo em identificar corretamente esse tipo de batimento. Esse comportamento é esperado, 
dado que essa classe apresenta menor representatividade e maior complexidade morfológica em comparação às demais.

\section{Resultados da CNN2D}

Nesta seção são apresentados os resultados obtidos pela CNN2D, utilizando como entrada os espectrogramas gerados a partir da Transformada de 
Fourier de Curta Duração (STFT).

A matriz de confusão correspondente pode ser vista na Figura \ref{fig:matconf2d}.

\begin{figure}[ht]
  \caption{Matriz de confusão da CNN2D}
  \label{fig:matconf2d}
  \centering
  \includegraphics[width=0.7\textwidth]{MatConf2D.jpeg}
  \legend{Elaborado pelo autor (2026)}
\end{figure}

A classe N apresentou confusões principalmente com as classes S, V e F, o que pode estar associado à similaridade entre padrões espectrais de batimentos 
normais e alguns tipos de arritmias, especialmente quando as variações são sutis no domínio da frequência. A classe S também apresentou confusão 
relevante com a classe N, indicando dificuldade do modelo em separar completamente essas categorias.

A classe V apresentou bom desempenho geral, com alta taxa de acerto, embora ainda ocorram algumas confusões com a classe F. Esse comportamento sugere 
que, mesmo no domínio espectral, certos padrões associados a batimentos ventriculares e de fusão permanecem próximos, dificultando sua distinção.

Um ponto crítico observado é novamente a confusão de batimentos da classe F com a classe N. Esse tipo de erro possui relevância clínica significativa, 
pois representa a classificação de batimentos anômalos como normais, o que pode levar à não detecção de eventos cardíacos importantes. Além disso, 
observa-se também confusão da classe F com a classe V, reforçando a complexidade na identificação desse tipo de batimento.

Novamente a classe Q apresentou bom desempenho, com altíssima taxa de acerto, indicando que seus padrões espectrais são bem definidos e facilmente 
separáveis pelo modelo.

Os valores calculados para as métricas de avaliação da CNN2D encontram-se na Tabela \ref{tab:table6}.

\begin{table}[ht]
  \caption{Métricas de desempenho da CNN2D}
  \label{tab:table6}
  \centering
  \resizebox{\textwidth}{!}{%
  \begin{tabular}{cccccc}
    \hline
    \bfseries Classe & \bfseries Acurácia & \bfseries Sensibilidade & \bfseries Especificidade & \bfseries Precisão & \bfseries F1-Score\\
    \hline
    N & 0.9570 & 0.9250 & 0.9650 & 0.8685 & 0.8959\\
    S & 0.9730 & 0.9000 & 0.9912 & 0.9626 & 0.9302\\
    V & 0.9720 & 0.9750 & 0.9712 & 0.8945 & 0.9330\\
    F & 0.9690 & 0.8800 & 0.9912 & 0.9617 & 0.9191\\
    Q & 0.9970 & 0.9900 & 0.9988 & 0.9950 & 0.9925\\
    \hline
    \textbf{MÉDIA} & \textbf{0.9736} & \textbf{0.9340} & \textbf{0.9835} & \textbf{0.9365} & \textbf{0.9341}\\
    \hline
  \end{tabular}%
  }
  \source
\end{table}

De maneira geral, o modelo apresentou desempenho elevado, com acurácia média de 97,36\% e F1-score médio de 93,41\%. Entretanto, observa-se que a 
sensibilidade da classe F (0,88) é inferior às demais, indicando maior dificuldade na identificação correta desse tipo de batimento. Esse resultado 
reforça a complexidade dessa classe e sugere que, mesmo com a utilização de representações espectro-temporais, ainda há limitações na discriminação de 
padrões mais sutis e menos representativos.

\section{Discussão e Comparativo}

Os resultados obtidos indicam que ambos os modelos propostos apresentaram bom desempenho na tarefa de classificação de batimentos cardíacos, com métricas 
médias bastante próximas entre si. A CNN1D apresentou desempenho ligeiramente superior em relação à CNN2D na maioria das métricas globais, especialmente 
em sensibilidade e F1-score. No entanto, essa diferença não é expressiva o suficiente para afirmar, de forma conclusiva, a superioridade de uma 
abordagem sobre a outra.

A análise das matrizes de confusão evidencia que ambos os modelos apresentam maior dificuldade na classificação das classes minoritárias, especialmente 
a classe F (batimentos de fusão). Em ambos os casos, observa-se uma tendência de confusão entre as classes F e N, o que pode ser explicado pela 
semelhança morfológica entre esses tipos de batimentos no domínio temporal. Esse comportamento é clinicamente relevante, pois a classificação incorreta 
de batimentos anômalos como normais pode comprometer a confiabilidade de sistemas automatizados de apoio ao diagnóstico.

As curvas de aprendizado dos modelos, apresentadas nas Figuras \ref{fig:curva1d} e \ref{fig:curva2d}, permitem uma análise mais aprofundada do comportamento durante o treinamento.

\begin{figure}[ht]
  \caption{Curvas CNN1D}
  \label{fig:curva1d}
  \centering
  \includegraphics[width=0.7\textwidth]{Curvas1D.png}
  \legend{Elaborado pelo autor (2026)}
\end{figure}

\begin{figure}[ht]
  \caption{Curvas CNN2D}
  \label{fig:curva2d}
  \centering
  \includegraphics[width=0.7\textwidth]{Curvas2D.png}
  \legend{Elaborado pelo autor (2026)}
\end{figure}

No caso da CNN1D, observa-se que a acurácia no conjunto de treinamento cresce rapidamente, aproximando-se de valores próximos a 100\%, enquanto a 
acurácia de validação estabiliza em torno de 94\%. De forma semelhante, a função de perda apresenta queda acentuada nas primeiras épocas, seguida por 
estabilização no conjunto de validação. Esse comportamento indica a presença de um leve sobreajuste, evidenciado pela diferença entre os desempenhos de 
treino e validação. No entanto, esse efeito é moderado e foi adequadamente controlado pelo uso de técnicas de regularização, como Dropout e Early 
Stopping, que interrompeu o treinamento por volta da época 14, evitando degradação do desempenho em validação.

Para a CNN2D, observa-se comportamento semelhante, porém com indícios mais pronunciados de sobreajuste. A acurácia no conjunto de treinamento atinge 
rapidamente valores próximos a 100\%, enquanto a acurácia de validação permanece em torno de 92–93\%, com maior instabilidade ao longo das épocas. 
Além disso, a função de perda no conjunto de validação apresenta oscilações mais evidentes após as primeiras épocas, sugerindo maior sensibilidade do 
modelo a variações nos dados de treinamento.

Comparativamente, os resultados indicam que, embora a CNN2D possua alta capacidade de aprendizado, a CNN1D apresentou melhor equilíbrio entre ajuste aos 
dados de treinamento e capacidade de generalização. Esse comportamento pode estar relacionado ao fato de que a representação temporal preserva 
diretamente as características morfológicas do sinal de ECG, enquanto a transformação para domínio espectro-temporal pode introduzir perda de informação 
ou simplificações decorrentes do redimensionamento das imagens.

De modo geral, as curvas de aprendizado evidenciam que ambos os modelos convergem rapidamente e se beneficiam do uso de mecanismos de regularização, 
sendo o critério de parada antecipada fundamental para evitar o sobreajuste e garantir melhor desempenho no conjunto de validação.

Quando comparados com trabalhos relacionados da literatura, os resultados obtidos neste estudo apresentam desempenho inferior. Essa diferença pode ser 
atribuída a diversos fatores. Primeiramente, muitos desses trabalhos utilizam arquiteturas mais profundas ou híbridas, combinando CNNs com outras 
abordagens, como redes recorrentes ou mecanismos de atenção, o que aumenta a capacidade de extração de características complexas.

Além disso, diferenças no pré-processamento dos sinais podem impactar significativamente o desempenho. Enquanto neste trabalho foram utilizados 
segmentos de 360 amostras centrados no batimento, outros estudos empregam janelas maiores ou técnicas adicionais de filtragem e remoção de ruído. Outro 
fator relevante é o protocolo de avaliação: alguns trabalhos utilizam validação cruzada entre pacientes (inter-patient), enquanto outros podem adotar 
divisões intra-paciente, o que tende a resultar em métricas mais elevadas devido à menor variabilidade entre os dados de treino e teste.

Outro aspecto importante diz respeito à definição dos hiperparâmetros. Neste trabalho, os parâmetros da rede, como número de filtros, tamanho dos 
kernels, taxa de dropout e taxa de aprendizado, foram definidos com base em práticas comuns da literatura e experimentações preliminares, sem a 
realização de uma busca sistemática. Essa escolha pode ter limitado o desempenho final dos modelos, uma vez que configurações mais otimizadas poderiam 
potencialmente produzir melhores resultados.

Adicionalmente, destaca-se que os resultados apresentados são provenientes de uma única execução para cada modelo, utilizando uma semente aleatória fixa. 
Dessa forma, não foi realizada uma análise estatística baseada em múltiplas execuções independentes, o que impede a avaliação da variabilidade dos 
resultados. Consequentemente, não é possível afirmar com rigor estatístico que as diferenças observadas entre os modelos são significativas, podendo 
estar dentro da variância natural do processo de treinamento.

Apesar dessas limitações, os resultados obtidos demonstram que tanto a abordagem baseada em sinais temporais (CNN1D) quanto a abordagem 
espectro-temporal (CNN2D) são capazes de capturar características relevantes dos sinais de ECG, apresentando desempenho competitivo para a tarefa proposta.